\section{Restoring the guarantee as a dial}
\label{sec:dial}

\begin{table}[t]
\centering
\footnotesize
\setlength{\tabcolsep}{3.5pt}
\begin{tabular}{llrrcrrrrr}
\toprule
Target & Controller & $\varepsilon$ & Commit viol.\ (gap) & 95\% CI & vs target & Breaches & Peak kVA & Bill Rs & Solve ms \\
\midrule
tight & marginal $q95$ & -- & \textbf{0.106} & [0.03, 0.20] & 0.049 & 4 & 459.0 & 1,138,193 & 6 \\
tight & scenario & 0.05 & \textbf{0.097} (+0.047) & [0.04, 0.17] & 0.097 & 9 & 464.6 & 1,141,532 & 52 \\
tight & scenario & 0.10 & \textbf{0.114} (+0.014) & [0.04, 0.21] & 0.096 & 9 & 466.8 & 1,142,186 & 52 \\
tight & scenario & 0.20 & \textbf{0.156} (-0.044) & [0.07, 0.25] & 0.073 & 10 & 469.2 & 1,142,647 & 55 \\
tight & scenario & 0.35 & \textbf{0.219} (-0.131) & [0.13, 0.33] & 0.077 & 13 & 467.5 & 1,141,036 & 55 \\
nominal & marginal $q95$ & -- & \textbf{0.103} & [0.03, 0.19] & 0.000 & 0 & 481.9 & 1,149,580 & 6 \\
nominal & scenario & 0.05 & \textbf{0.078} (+0.028) & [0.02, 0.15] & 0.000 & 0 & 480.5 & 1,148,576 & 63 \\
nominal & scenario & 0.10 & \textbf{0.097} (-0.003) & [0.03, 0.18] & 0.000 & 0 & 486.0 & 1,151,865 & 63 \\
nominal & scenario & 0.20 & \textbf{0.150} (-0.050) & [0.07, 0.24] & 0.000 & 0 & 480.1 & 1,148,071 & 63 \\
nominal & scenario & 0.35 & \textbf{0.208} (-0.142) & [0.12, 0.31] & 0.000 & 0 & 484.0 & 1,150,364 & 55 \\
\bottomrule
\end{tabular}
\caption{Closed loop over one billing month on a 13{,}759\,m$^2$ office in Phoenix under the TNERC tariff. \textbf{Commit violation} is the acceptance metric: the fraction of horizons in which realised load cleared the ceiling the optimiser committed to, which is what $\varepsilon$ is a statement about. The against-target column is the business metric and is not what the chance constraint promises, since the committed peak is a decision variable. Rank correlation 0.957, mean absolute gap 0.057, conservative at 5 of 8 levels; resolution floor $1/S=0.025$. Intervals are a day-block bootstrap over the sampled commitments, which are overlapping 16-hour windows and so are not the independent draws a binomial interval would assume; the requested $\varepsilon$ lies inside its own interval at 6 of 8 levels, and the mean half-width is $0.081$, which is the resolution one billing month buys.}
\label{tab:acceptance}
\end{table}


We fit a Student-$t$ copula over the horizon and sample joint trajectories
--- the marginals-to-scenarios construction of \citet{pinson2009scenarios},
with a heavier-tailed dependence family --- which a scenario MILP then
constrains with a CVaR surrogate \citep{rockafellar2000cvar,nemirovski2006convex}
on the probability mass permitted to violate. The copula reproduces the calibrated marginals to
within $4.6\%$ of mean interval width at $4{,}000$ paths per origin, so it adds
dependence and disturbs nothing upstream, and it predicts horizon exceedance far
closer than either substitute: $0.497$ against a realised $0.405$ (interval
$0.28$--$0.53$) at $H=64$,
where \eqref{eq:marginal} implies $0.05$ and the independence figure is $0.983$.

Table~\ref{tab:acceptance} is a closed-loop run over a real billing month at two
demand targets and four risk levels. The acceptance metric is violation of the
ceiling \emph{the optimiser itself committed to at each origin}, because that is
the event $\varepsilon$ is a statement about. The operator's monthly demand
target is a business metric the chance constraint never promised: the committed
peak is a decision variable, and an optimiser short of capacity may lift it, pay
the demand charge, and honour its constraint exactly. Grading a chance
constraint against a number it never promised is the usual way this test is made
meaningless, so both columns are shown and only one is the acceptance criterion.

The result is that $\varepsilon$ becomes a dial. Realised commit violation
tracks the requested level with rank correlation $0.957$ across the sweep, mean
absolute gap $0.057$, and lands conservative at five of eight levels. The
formulation errs safe by construction: the CVaR surrogate admits a strict subset
of the chance-constrained feasible set, so it defends the committed ceiling
harder than asked. Those are point estimates over one month, and the
commitments behind them are overlapping sixteen-hour windows, so each rate
carries a day-block interval; the mean half-width is $0.081$. The requested
$\varepsilon$ lies inside its own interval at six of eight levels, and the two
it does not are both $\varepsilon=0.35$, on the safe side: on this plant the
dial is never measurably unsafe.

\begin{table}[t]
\centering
\footnotesize
\setlength{\tabcolsep}{3.5pt}
\begin{tabular}{llrrcrrrrr}
\toprule
Target & Controller & $\varepsilon$ & Commit viol.\ (gap) & 95\% CI & vs target & Breaches & Peak kVA & Bill Rs & Solve ms \\
\midrule
tight & marginal $q95$ & -- & \textbf{0.297} & [0.18, 0.42] & 0.101 & 13 & 890.9 & 2,837,240 & 6 \\
tight & scenario & 0.05 & \textbf{0.178} (+0.128) & [0.08, 0.28] & 0.101 & 14 & 885.5 & 2,833,912 & 51 \\
tight & scenario & 0.10 & \textbf{0.233} (+0.133) & [0.13, 0.35] & 0.124 & 15 & 895.3 & 2,835,693 & 61 \\
tight & scenario & 0.20 & \textbf{0.300} (+0.100) & [0.19, 0.42] & 0.123 & 14 & 903.8 & 2,838,637 & 62 \\
tight & scenario & 0.35 & \textbf{0.392} (+0.042) & [0.26, 0.52] & 0.075 & 9 & 915.4 & 2,842,287 & 63 \\
nominal & marginal $q95$ & -- & \textbf{0.297} & [0.18, 0.42] & 0.000 & 0 & 914.0 & 2,843,724 & 6 \\
nominal & scenario & 0.05 & \textbf{0.172} (+0.122) & [0.08, 0.27] & 0.023 & 1 & 945.5 & 2,849,643 & 61 \\
nominal & scenario & 0.10 & \textbf{0.236} (+0.136) & [0.13, 0.36] & 0.023 & 1 & 932.7 & 2,846,024 & 63 \\
nominal & scenario & 0.20 & \textbf{0.300} (+0.100) & [0.19, 0.42] & 0.023 & 1 & 931.3 & 2,846,298 & 81 \\
nominal & scenario & 0.35 & \textbf{0.392} (+0.042) & [0.26, 0.52] & 0.023 & 1 & 938.9 & 2,848,913 & 59 \\
\bottomrule
\end{tabular}
\caption{Closed loop over one billing month on the IIIT-Delhi campus feed under the DERC tariff, the billed consumer of Section~\ref{sec:data}. \textbf{Commit violation} is the acceptance metric: the fraction of horizons in which realised load cleared the ceiling the optimiser committed to, which is what $\varepsilon$ is a statement about. The against-target column is the business metric and is not what the chance constraint promises, since the committed peak is a decision variable. Rank correlation 0.988, mean absolute gap 0.100, conservative at 0 of 8 levels; resolution floor $1/S=0.025$. Intervals are a day-block bootstrap over the sampled commitments, which are overlapping 16-hour windows and so are not the independent draws a binomial interval would assume; the requested $\varepsilon$ lies inside its own interval at 4 of 8 levels, and the mean half-width is $0.113$, which is the resolution one billing month buys.}
\label{tab:acceptance-campus}
\end{table}


\paragraph{A second control object, and the half of this that replicates.}
Table~\ref{tab:acceptance-campus} repeats the sweep on the IIIT-Delhi campus
feed --- another country, another tariff, another plant, and the consumer a
Delhi utility actually bills --- with nothing tuned for it. Two things happen,
and only one of them is the result above.

The dial's \emph{monotonicity} replicates, and more cleanly than on the office:
rank correlation $0.988$, and realised commit violation rises through $0.172$,
$0.236$, $0.300$, $0.392$ as $\varepsilon$ is set to $0.05$, $0.10$, $0.20$,
$0.35$. An operator turning the dial gets a monotone response on both plants.

The dial's \emph{calibration} does not replicate. Mean absolute gap is $0.100$
against $0.057$, and the sweep is conservative at \emph{none} of its eight
levels rather than five: on this feed the formulation delivers more risk than
it was asked for at every setting. The intervals sharpen where that matters:
they exclude the requested level at $\varepsilon=0.05$ and $0.10$ on both
targets --- the two tightest settings, the ones an operator would actually use
--- and contain it at $0.20$ and $0.35$, four of eight in all. We do not think this is the risk
formulation failing, and the reason is one row of
Table~\ref{tab:horizon-panel}. The campus feed is the single window in the
study whose per-step layer \emph{under}-covers --- $\hat\alpha = 0.079$
against a nominal $0.05$, the only figure above the line --- and $\varepsilon$
is a statement built on top of that layer. A joint guarantee cannot be more
exact than the marginals it composes, and here the marginals are optimistic by
half again.

That is a testable claim, so we tested it. Restating the adaptive step in the
interval's units --- which Section~\ref{sec:calibration} shows is the
defensible way to set it --- moves the campus per-step rate from $0.079$ to
$0.069$, and the whole sweep was rerun on that calibration with nothing else
changed. The dial's exactness improves in the direction the explanation
predicts, and by an amount this experiment cannot separate from noise. Mean
absolute gap falls from $0.100$ to $0.083$ and realised violation at
$\varepsilon=0.05$ from $0.178$ to $0.158$, both toward the $0.057$ and $0.097$
of the well-calibrated plant; rank correlation is unchanged at $0.988$ and the
sweep is still conservative at none of its eight levels. But each of those
rates is a fraction of the commitments sampled over one month, and the
day-block interval on it is wider than the movement --- a mean half-width of
$0.113$ on the campus sweep against a shift of $0.018$ --- so what this run
establishes is a direction and not a size. The claim the paper makes is
therefore the weak one: a per-step layer that under-covers is a plausible and
partially-supported cause of a dial that over-delivers risk, the repair moves
the dial the right way, and the remaining distance to the calibrated plant is
neither shown to close nor shown to persist. Separating those would need more
billing months than one.

The honest summary of Contribution~2 is therefore narrower than one control
object suggested. On both plants $\varepsilon$ converts a risk level that was
an emergent property of horizon length into one an operator sets and gets a
monotone response from. On the plant whose marginals are calibrated it is also
roughly exact and errs safe; on the plant whose marginals are not, it is
monotone and optimistic, and repairing those marginals moves it the right way
by an amount one month cannot resolve. That is a useful thing for a deployment to know, and
it says the marginal calibration layer is not a preliminary to the guarantee
but part of it.

\paragraph{What it does not do, stated in full.} It is not exact. At the
tightest setting, $\varepsilon=0.05$ delivers $0.097$ at the tight target,
nearly twice the requested risk and on the unsafe side --- though one month's
interval on it, $0.04$--$0.17$, cannot tell that miss from none. The dial has
two resolution limits. One is statistical: the intervals are wider than the
spacing between adjacent levels, so the monotone ordering above is a property
of the point estimates, and resolving neighbouring settings one from another
would take more billing months, not more scenarios. The other is a hard floor
at $1/S = 0.025$: with $S$ scenarios the finest expressible
risk level is one scenario's worth of probability mass, below which the CVaR
surrogate collapses onto the sample maximum and successive $\varepsilon$ cease
to be distinguishable. That is a property of sample average approximation rather
than a tuning artefact, and a finer dial costs scenarios and solve time.

\begin{table}[t]
\centering
\footnotesize
\setlength{\tabcolsep}{3.5pt}
\begin{tabular}{llrcrrrrr}
\toprule
Series & June & realised & 95\% CI & $\nu_{\mathrm{tail}}$ & copula & $\nu_{\mathrm{lik}}$ & copula & copula, $\nu{=}7$ \\
\midrule
Phoenix office (BDG2) & 2017 & \textbf{0.405} & [0.28, 0.53] & 7 & 0.497 & 25 & 0.567$^\dagger$ & 0.497 \\
\midrule
IIIT-Delhi Academic & 2015 & \textbf{0.393} & [0.28, 0.50] & 3 & 0.477 & 25 & 0.714$^\dagger$ & 0.616$^\dagger$ \\
IIIT-Delhi Academic & 2016 & \textbf{0.636} & [0.50, 0.75] & 5 & 0.580 & 25 & 0.732 & 0.631 \\
IIIT-Delhi Academic & 2017 & \textbf{0.422} & [0.30, 0.54] & 5 & 0.595$^\dagger$ & Gauss. & 0.801$^\dagger$ & 0.648$^\dagger$ \\
IIIT-Delhi Boys & 2014 & \textbf{0.501} & [0.38, 0.61] & 4 & 0.554 & 25 & 0.746$^\dagger$ & 0.642$^\dagger$ \\
IIIT-Delhi Campus & 2017 & \textbf{0.432} & [0.29, 0.58] & 4 & 0.482 & 25 & 0.636$^\dagger$ & 0.553 \\
IIIT-Delhi Dining & 2017 & \textbf{0.714} & [0.61, 0.81] & 15 & 0.744 & 25 & 0.774 & 0.665 \\
IIIT-Delhi Facilities & 2015 & \textbf{0.639} & [0.55, 0.71] & 7 & 0.705 & Gauss. & 0.875$^\dagger$ & 0.705 \\
IIIT-Delhi Facilities & 2016 & \textbf{0.621} & [0.52, 0.72] & 5 & 0.633 & Gauss. & 0.856$^\dagger$ & 0.690 \\
IIIT-Delhi Facilities & 2017 & \textbf{0.739} & [0.62, 0.85] & 10 & 0.753 & Gauss. & 0.876$^\dagger$ & 0.706 \\
IIIT-Delhi Girls & 2014 & \textbf{0.693} & [0.59, 0.78] & 4 & 0.589$^\dagger$ & Gauss. & 0.848$^\dagger$ & 0.684 \\
IIIT-Delhi Girls & 2017 & \textbf{0.818} & [0.74, 0.89] & 7 & 0.739 & Gauss. & 0.914$^\dagger$ & 0.739 \\
IIIT-Delhi Library & 2014 & \textbf{0.480} & [0.36, 0.59] & 4 & 0.567 & Gauss. & 0.816$^\dagger$ & 0.659$^\dagger$ \\
IIIT-Delhi Library & 2016 & \textbf{0.509} & [0.41, 0.61] & 4 & 0.504 & 25 & 0.670$^\dagger$ & 0.579 \\
\bottomrule
\end{tabular}
\caption{What selecting the copula's degrees of freedom contributes. Realised horizon exceedance at $H=64$ on the held-out month with its day-block bootstrap interval, and the copula's prediction under three choices of $\nu$: matched to horizon exceedance on the validation block (the paper's choice; columns 5--6), maximising the $t$-copula pseudo-likelihood on the same block (columns 7--8), and held at the Phoenix value $\nu=7$ with no refitting (last column). $\dagger$ marks a prediction outside the realised interval. The tail-matched choice lies inside on 12 of 14 windows (mean absolute error $0.065$). The likelihood criterion selects $\nu \geq 25$ or the Gaussian on every window, over-predicts on 14 of 14 and lies inside on 2: it is dominated by the body of the distribution, where the families agree, and cannot see the upper-tail dependence the horizon event is made of. The Phoenix value applied unchanged to the 13 Indian windows lies inside on 9 of 13 (mean absolute error $0.098$); the per-window selection buys accuracy, not the agreement itself. The independence figure lies inside on none.}
\label{tab:copula-df}
\end{table}


Two further caveats belong here rather than in a footnote, and one that was a
caveat has been measured. The copula's degrees of freedom are selected on the
validation block by matching predicted to realised horizon exceedance, which is
the quantity the copula is later judged on; the split is honest, but the
criterion knows what it will be graded on. Table~\ref{tab:copula-df} bounds what
that choice contributes. Holding $\nu$ at the Phoenix value of $7$ and applying
it unchanged to the thirteen Indian windows, the copula still lies inside the
realised interval on nine of thirteen, with mean absolute error $0.098$ against
$0.063$ when $\nu$ is selected per window: the selection buys accuracy, not the
agreement. The textbook alternative --- choosing $\nu$ by the $t$-copula
pseudo-likelihood on the same validation block --- selects $\nu \geq 25$ or the
Gaussian on every one of the fourteen windows, over-predicts on all fourteen,
and lies inside the interval on two. The likelihood is dominated by the body of
the distribution, where the two families agree; the horizon event is a statement
about the joint upper tail, where they do not. That is why the one free
parameter is fitted to the decision-relevant statistic rather than to the
likelihood, and it is a warning for anyone fitting a dependence model for a
tail event by the default criterion. Base-load and PV scenarios are drawn
independently, while cloud physically raises cooling load at the same moment it
cuts PV output, so the true joint law has worse afternoons than these scenarios
contain. And the horizon is sixteen hours while the demand charge bills a
monthly maximum over roughly $2{,}880$ intervals: a joint guarantee over the
window is strictly stronger than a marginal one and strictly weaker than a
monthly one. Extending it needs the copula fitted over a month-length horizon,
which is a $2{,}880$-square correlation matrix and needs structure rather than a
raw estimate.

The honest summary is not that the guarantee is restored exactly. It is that a
risk level which was previously an \emph{emergent property of horizon length}
--- unstated, unsettable, and equal to $0.106$ at the tight target on the
office and $0.297$ on the campus feed under the incumbent marginal controller
--- is now a parameter with a monotone response on both, conservative on the
one whose marginals are calibrated, at a cost of \rupee{3{,}339} on the month
there and \rupee{3{,}329} \emph{saved} on the campus, where the tight target
binds so rarely that the scenario controller buys its risk level out of
headroom the marginal one was wasting.
